\section{模型的评价与拓展应用}

\subsection{模型的特色与优势}

本文针对复杂山地洪涝灾害应急场景，构建了多旋翼无人机应急运输与通信协同优化模型。模型具有以下主要特点：

\begin{enumerate}
    \item \textbf{融合高程信息与非线性动力学机理}：
    结合 30 米 DEM 高程数据，模型对爬升克服重力做功与水平飞行气动阻力进行分段计算；基于 $(3/2)$ 次方等效航程规律，利用二分搜索确定各机型在不同服务区的最大安全载荷边界，提高了复杂起伏地形下的飞行时耗与电量消耗估算精度。
    
    \item \textbf{机体与动力电池解耦排程}：
    针对多旋翼无人机充电时间较长、单次航时受限的特征，建立了机体与共享电池解耦排程模型。引入两阶段非线性等效充电函数，通过离散事件仿真统筹 8 架实体机与 14 组共享电池的流转过程，避免充放电时序冲突，将系统最大完工时间控制在 2.15 小时（7730.5 s）。
    
    \item \textbf{射线追踪与第一菲涅尔区视距判定}：
    在空地通信建模中，利用三维空间射线追踪判断山体对视距通信的遮挡情况；结合第一菲涅尔区几何要求筛选中继机悬停制高点，以 3 个中继架次实现了运输机全航程连续视距通信。
    
    \item \textbf{图论连通分支与分区调度分析}：
    分析了多点运输回路强绑定约束下的等价关系，并利用图论连通分支给出了 $K=2$ 与 $K=3$ 任务分区的可行分解方案；定量测算了独立封闭分区管理对无人机和电池资源复用造成的约束与设备缺口，为应急管理提供了理论依据。
\end{enumerate}

\subsection{模型的模块化设计与拓展能力}

本文模型采用松耦合的模块化设计，具备良好的扩展能力：

\begin{enumerate}
    \item \textbf{三维动态风场接口扩展}：
    动力学模块在物理计算中将重力做功与气动阻力分段解耦。在实际应用中，可直接将中尺度气象预报（WRF）或局部计算流体力学（CFD）风场数据代入各航段相对空速计算，评估不同风向与风速对飞行能耗的影响。
    
    \item \textbf{复杂电磁信道模型扩展}：
    当前射线追踪与菲涅尔区几何判断已建立了三维空间遮挡骨架。未来可进一步结合地表介电常数与植被分布，引入抛物线方程（PE）等高阶电磁传播算法或多径衰落信道模型，以适应更复杂的深谷通信场景。
    
    \item \textbf{大规模调度快速求解拓展}：
    本文采用的图论宏观定界与启发式搜索相分离的求解结构，有利于控制计算复杂度。离线仿真生成的全流程调度时序数据，亦可作为样本库用于训练图神经网络或强化学习策略，以提升突发场景下的在线响应速度。
\end{enumerate}

\subsection{实际应用建议}

根据本文的模型推导与计算结果，针对实际应急救援行动提出以下建议：

\begin{enumerate}
    \item \textbf{推广地理信息与动力学校核一体化的航线规划}：
    建议在应急调度系统中集成高精度 DEM 数据与无人机动力学校核算法，在灾情初期快速自动生成净空安全航线、各航段最大有效载荷以及视距盲区分布，减少人工经验估算带来的航程不足与撞山风险。
    
    \item \textbf{配置标准化移动充换电装备以提高无人机利用率}：
    计算表明，合理配置备用电池与充电流水线可有效减少机体等待时间。建议救援单位按照约 1:2 至 1:3 的机电比配置备用电池，并在前进基地配备快充设备，缩短机体地面等待时间，保证作业连续性。
    
    \item \textbf{统筹分区责任制与跨区机动调度机制}：
    第四问计算表明，各区域完全独立封闭会导致设备无法跨组复用，从而产生明显的设备缺口。建议在落实属地责任的同时，允许中继机或高负荷运输机在空闲窗口接受统一机动调度；在装备储备方面，适当增加主力机型（如 B 型机）与备用电池的储备冗余。
\end{enumerate}
